% ============================================================
%  SCHOLA DOMESTICA — Magister Mathematica
%  Level 1 | Week 10 | Day 3
%  Topic: Near-Doubles Strategy & Mixed Facts: Sums to 8
% ============================================================
\documentclass[12pt, letterpaper]{article}

\usepackage[top=1in, bottom=1in, left=1in, right=1in]{geometry}
\usepackage[T1]{fontenc}
\usepackage[utf8]{inputenc}
\usepackage{lmodern}
\usepackage{microtype}
\usepackage{amsmath}
\usepackage{xcolor}
\usepackage{tikz}
\usepackage{array}
\usepackage{enumitem}
\usepackage{multicol}
\usepackage{mdframed}
\usepackage{fancyhdr}

\definecolor{scholared}{RGB}{160,30,30}
\definecolor{scholagold}{RGB}{180,140,40}
\definecolor{scholacream}{RGB}{253,248,235}
\definecolor{scholagreen}{RGB}{50,110,60}
\definecolor{scholablue}{RGB}{40,80,140}

\pagestyle{fancy}
\fancyhf{}
\renewcommand{\headrulewidth}{0.6pt}
\fancyhead[L]{\small\textcolor{scholared}{\textbf{Schola Domestica — Magister Mathematica}}}
\fancyhead[R]{\small\textcolor{scholared}{Level 1 · Week 10 · Day 3}}
\fancyfoot[C]{\small\textcolor{gray}{— \thepage\ —}}

\newcommand{\answerline}{\underline{\hspace{2.5cm}}}
\newcommand{\biganswerline}{\underline{\hspace{4cm}}}
\newcommand{\numbox}{\tikz{\draw[scholablue,thick](0,0)rectangle(1.2cm,0.85cm);}}

\newmdenv[
  backgroundcolor=scholacream,
  linecolor=scholared,
  linewidth=1.5pt,
  roundcorner=6pt,
  innertopmargin=8pt,
  innerbottommargin=8pt,
  innerleftmargin=10pt,
  innerrightmargin=10pt
]{lessonbox}

\newmdenv[
  backgroundcolor={rgb,255:red,235;green,245;blue,235},
  linecolor=scholagreen,
  linewidth=1.2pt,
  roundcorner=4pt,
  innertopmargin=6pt,
  innerbottommargin=6pt,
  innerleftmargin=8pt,
  innerrightmargin=8pt
]{enrichbox}

\begin{document}

\begin{center}
  {\LARGE\textbf{\textcolor{scholared}{Near-Doubles:}}}\\[4pt]
  {\Large\textcolor{scholagold}{\textit{Almost the Same — Just One More!}}}\\[6pt]
  {\small Name: \biganswerline \hspace{1cm} Date: \answerline}
\end{center}

\vspace{4pt}
\noindent\textcolor{scholared}{\rule{\linewidth}{1pt}}
\vspace{4pt}

% IMAGE PROMPT (Beatrix Potter watercolour style):
% Two robin chicks side by side in a nest; one has an extra worm dangling from its beak. The mother
% robin watches from a branch. Warm brown and red-breast palette, soft watercolour, no text.

\begin{lessonbox}
\textbf{Near-doubles} are facts that are \emph{almost} a double — one part is one more.\\[6pt]
\noindent
\begin{tabular}{@{}p{6cm} p{6cm}@{}}
\textbf{Know the double:} & \textbf{Use it for the near-double:} \\[4pt]
$3 + 3 = 6$ & $3 + 4 = 7$ \quad (one more than 6) \\[3pt]
$4 + 4 = 8$ & $4 + 3 = 7$ \quad (one less than 8)
\end{tabular}
\medskip

\noindent\textit{Think: double, then add 1 or take away 1.}
\end{lessonbox}

\bigskip

% ===== PART I: IDENTIFY DOUBLE AND NEAR-DOUBLE =====
\noindent\textcolor{scholared}{\large\textbf{Part I — Double or Near-Double?}} \hfill \textit{(Problems 1–6)}

\medskip
\noindent\textit{Write D if it is a double, N if it is a near-double. Then solve.}

\bigskip

\noindent
\begin{tabular}{@{}p{3.6cm} p{3.6cm} p{3.6cm} p{3.6cm}@{}}
\textbf{1.}\; $3 + 3 =$ \numbox \; \underline{\hspace{0.6cm}} &
\textbf{2.}\; $3 + 4 =$ \numbox \; \underline{\hspace{0.6cm}} &
\textbf{3.}\; $2 + 2 =$ \numbox \; \underline{\hspace{0.6cm}} &
\textbf{4.}\; $2 + 3 =$ \numbox \; \underline{\hspace{0.6cm}} \\[14pt]
\multicolumn{2}{@{}l@{}}{\textbf{5.}\; $4 + 4 =$ \numbox \; \underline{\hspace{0.6cm}}} &
\multicolumn{2}{@{}l@{}}{\textbf{6.}\; $4 + 5 =$ \numbox \; \underline{\hspace{0.6cm}}}
\end{tabular}

\bigskip\bigskip

% ===== PART II: USE THE DOUBLE TO FIND THE NEAR-DOUBLE =====
\noindent\textcolor{scholared}{\large\textbf{Part II — Step-by-Step Near-Doubles}} \hfill \textit{(Problems 7–10)}

\medskip
\noindent\textit{Fill in both steps.}

\bigskip

\noindent\textbf{7.} Double: $3 + 3 =$ \numbox \hspace{1cm} Near-double: $3 + 4 =$ \numbox

\bigskip

\noindent\textbf{8.} Double: $4 + 4 =$ \numbox \hspace{1cm} Near-double: $4 + 3 =$ \numbox

\bigskip

\noindent\textbf{9.} Double: $2 + 2 =$ \numbox \hspace{1cm} Near-double: $2 + 3 =$ \numbox

\bigskip

\noindent\textbf{10.} Double: $1 + 1 =$ \numbox \hspace{1cm} Near-double: $1 + 2 =$ \numbox

\bigskip\bigskip

% ===== PART III: MIXED FACTS SUMS TO 8 =====
\noindent\textcolor{scholared}{\large\textbf{Part III — Mixed Facts: All Sums to 8}} \hfill \textit{(Problems 11–16)}

\medskip

\noindent
\begin{tabular}{@{}p{3.3cm} p{3.3cm} p{3.3cm} p{3.3cm}@{}}
\textbf{11.}\; $5 + 3 =$ \numbox &
\textbf{12.}\; $4 + 4 =$ \numbox &
\textbf{13.}\; $6 + 2 =$ \numbox &
\textbf{14.}\; $7 + 1 =$ \numbox \\[12pt]
\multicolumn{2}{@{}l@{}}{\textbf{15.}\; $5 + 2 =$ \numbox} &
\multicolumn{2}{@{}l@{}}{\textbf{16.}\; $3 + 4 =$ \numbox}
\end{tabular}

\bigskip

\begin{enrichbox}
\textbf{\textcolor{scholagreen}{Near-Doubles Chart}}\\[4pt]
\noindent
\begin{tabular}{|c|c|c|}
\hline
\rule{0pt}{11pt}\textbf{Double} & \textbf{Sum} & \textbf{Near-double} \\
\hline
$1+1$ & 2 & $1+2=3$ \\
$2+2$ & 4 & $2+3=5$ \\
$3+3$ & 6 & $3+4=7$ \\
$4+4$ & 8 & $4+5=9$ \\
\hline
\end{tabular}

\medskip
\noindent\textit{Doubles are anchors. Near-doubles are just one step away!}
\end{enrichbox}

\vspace{0.3cm}
\noindent\textcolor{gray}{\small\textit{Ray's Primary Arithmetic — using known facts to find new ones.}}

\end{document}
